% Phase-Toroidal Engine — arXiv submission
\documentclass[12pt,a4paper]{article}

\usepackage[utf8]{inputenc}
\usepackage[T1]{fontenc}
\usepackage{amsmath,amssymb,amsfonts}
\usepackage{graphicx}
\usepackage{hyperref}
\usepackage{cite}
\usepackage{geometry}
\usepackage{abstract}

\geometry{margin=2.5cm}

\newcommand{\arXiv}[1]{\href{https://arxiv.org/abs/#1}{arXiv:#1}}

\title{Phase-Toroidal Engine: \\
A Theoretical Framework for Vacuum-Mediated Propulsion \\
via Coupled Topological Phase Transitions}

\author{IAC Research Group \\
\texttt{https://t.me/my\_super\_trader\_v32\_bot} \\
IP Protection: Zenodo 10.5281/zenodo.21209447}

\date{\today}

\begin{document}

\maketitle

\begin{abstract}
We present a theoretical framework for a novel class of propellantless propulsion -- the Phase-Toroidal Engine (PTE) -- which generates directed impulse by coupling toroidal field topology to controlled vacuum phase transitions. The engine exploits three coupled mechanisms: (A) nested toroidal electromagnetic confinement producing a symmetry-breaking stress tensor along the torus axis; (B) scalar dilaton mode coupling that locally modulates the vacuum refractive index, creating a metric gradient; and (C) controlled cycling through a critical topological phase transition in the vacuum state, yielding net impulse during the symmetry-breaking phase. We derive the governing effective Lagrangian, estimate thrust-to-power scaling, identify parametric stability conditions, and propose concrete experimental validation pathways. The framework naturally connects to recent empirical results: the 95.4 GeV dilaton resonance (CMS/ATLAS), toroidal spin ice Blume-Capel tricritical phase transitions, and topological vacuum energy modulation. Six falsifiable predictions are identified, spanning Casimir force anomalies at golden-ratio resonance spacings, anomalous thrust from asymmetric RF cavities, neutron star spin-down anomalies, and gamma-ray signatures from vacuum phase transition decay. The PTE framework offers a testable bridge between quantum vacuum physics, condensed matter topology, and practical aerospace engineering.
\end{abstract}

\section{Introduction}

The fundamental challenge of space propulsion remains the tyranny of the rocket equation: every kilogram of reaction mass must be accelerated along with the payload, creating an exponential mass penalty for delta-v. Propellantless propulsion -- generating thrust without onboard reaction mass -- has been a central aspiration of aerospace physics.

Recent decades have produced serious theoretical proposals, including the Woodward effect (transient mass fluctuations from Mach's principle), the Alcubierre metric requiring exotic matter with negative energy density, and asymmetric RF cavity drives. None have yet produced a theoretically consistent and experimentally reproducible framework grounded in conservative physics.

The Phase-Toroidal Engine (PTE) proposed here takes a different approach. Rather than invoking exotic matter, it exploits topological phase transitions in the quantum vacuum -- a physically well-motivated phenomenon -- coupled to toroidal field geometry to produce directed impulse. The central insight is that the vacuum can undergo phase transitions under external fields, and that symmetry-breaking dynamics can be geometrically directed.

This work integrates several recent theoretical developments:
\begin{itemize}
    \item The Toroidal Phase-Transition Model (TPTM) replaces gravitational singularities with rotating torus configurations~\cite{weinstock2026}
    \item The Photonic-Conjugated Model (PCM) models massive particles as confined light in toroidal topologies~\cite{caputo2025}
    \item The Toroidal Torsion Field Theory (TTFT) derives fundamental constants from golden-ratio eigenmode hierarchies on toroidal manifolds~\cite{zero2025}
    \item The 95.4 GeV dilaton resonance detected by CMS and ATLAS provides evidence for a scalar channel coupling to the trace anomaly~\cite{hofseth2026}
    \item Blume-Capel tricritical phase transitions have been realized in toroidal spin ice arrays~\cite{blume2025}
\end{itemize}

\section{Theoretical Framework}

\subsection{Toroidal Field Topology}

Consider a toroidal volume with major radius $R$ and minor radius $r$, parameterized by toroidal angle $\phi \in [0, 2\pi)$ and poloidal angle $\theta \in [0, 2\pi)$. The metric in toroidal coordinates is:

\begin{equation}
ds^2 = -dt^2 + (R + r\cos\theta)^2 d\phi^2 + r^2 d\theta^2 + dr^2
\label{eq:metric}
\end{equation}

Within this volume, we define a nested electromagnetic field configuration comprising a toroidal component $B_\phi$ winding the long way and a poloidal component $B_\theta$ winding the short way:

\begin{align}
B_\phi &= B_0 f(r/R) \\
B_\theta &= \varepsilon B_0 g(r/R)
\end{align}

where $\varepsilon$ is the asymmetry parameter and $f,g$ are radial profile functions satisfying the Grad-Shafranov equilibrium:

\begin{equation}
\Delta^* \psi = -\mu_0 R^2 \frac{dP}{d\psi} - F\frac{dF}{d\psi}
\label{eq:gradshafranov}
\end{equation}

Here $\psi$ is the poloidal flux function, $P$ the plasma pressure, and $F$ the toroidal field function.

The key geometric property for propulsion is the \textbf{central null zone} -- a region near the torus axis where the poloidal gradient of the combined field stress tensor is maximal while the total field pressure gradient vanishes:

\begin{align}
\nabla \cdot \mathbf{T} &\neq 0 \quad \text{(along symmetry axis)} \\
\nabla P_{\text{total}} &= 0 \quad \text{(in null zone center)}
\end{align}

This creates an uncompensated stress along the torus symmetry axis -- the foundation for directed impulse.

\subsection{Vacuum Phase Transitions}

Following the Refractive Vacuum Gravity (RVG) formalism~\cite{hofseth2026}, we model the vacuum as a medium with refractive index $K$ that responds to scalar fields coupled to the trace of the energy-momentum tensor:

\begin{equation}
K(x) = K_0 \exp(\alpha \phi(x) / M_{\text{Pl}})
\label{eq:refractive}
\end{equation}

where $\phi$ is a dilaton-like scalar field, $\alpha$ is a coupling constant, and $M_{\text{Pl}}$ is the Planck mass.

The vacuum effective action includes a phase transition term:

\begin{equation}
S_{\text{vac}} = \int d^4x \sqrt{-g} \left[ \frac{1}{2}(\partial\phi)^2 - V(\phi) + \frac{\phi}{4\Lambda} F_{\mu\nu}F^{\mu\nu} \right]
\label{eq:action}
\end{equation}

The potential $V(\phi)$ is assumed to have a tricritical point structure analogous to the Blume-Capel model~\cite{blume2025}:

\begin{equation}
V(\phi) = a_2(T)\phi^2 + a_4\phi^4 + a_6\phi^6
\label{eq:potential}
\end{equation}

where $a_2(T)$ changes sign at a critical temperature-like parameter, producing a first-order phase transition below a tricritical point and second-order above it.

Biasing the system through an asymmetric toroidal field shifts the effective potential:

\begin{equation}
V_{\text{eff}}(\phi) = V(\phi) - \frac{1}{2}\xi \mathbf{B}^2\phi
\label{eq:biased}
\end{equation}

where $\xi$ is the field coupling coefficient. The sign and magnitude of $\xi$ depend on the field topology. The toroidal asymmetry $\varepsilon$ creates a directional bias in the phase transition dynamics.

\subsection{Coupled Dynamics}

The core mechanism couples toroidal field topology to the vacuum phase transition through three phases:

\begin{enumerate}
    \item \textbf{Pumping phase:} Nested toroidal fields reach critical energy density $U_{\text{crit}} = \xi^{-1} V'(\phi_0)$
    \item \textbf{Transition phase:} Local vacuum undergoes topological phase transition (symmetry breaking), releasing energy $\Delta E = V(\phi_{\text{sym}}) - V(\phi_{\text{broken}})$
    \item \textbf{Impulse phase:} The asymmetry $\varepsilon$ directs the released energy along the torus symmetry axis
\end{enumerate}

The effective Lagrangian for the coupled system is:

\begin{align}
\mathcal{L}_{\text{PTE}} = &-\frac{1}{4}F_{\mu\nu}F^{\mu\nu} + \frac{1}{2}(\partial\phi)^2 - V(\phi) \nonumber \\
&+ \frac{\phi}{4\Lambda}F_{\mu\nu}F^{\mu\nu} + \epsilon^{\mu\nu\rho\sigma}A_\mu\partial_\nu\phi\partial_\rho A_\sigma
\label{eq:lagrangian}
\end{align}

The last term is a topological coupling that -- for toroidal boundary conditions -- produces a net axial current:

\begin{equation}
J_{\text{axial}} = \int d^3x \, \nabla \times (\phi \mathbf{B}) \neq 0
\label{eq:axial}
\end{equation}

\section{Thrust Estimate}

From the energy-momentum tensor of the coupled system, the thrust per cycle is:

\begin{equation}
F_{\text{cycle}} \approx \eta \cdot \frac{\Delta V}{c} \cdot \frac{R}{r} \cdot \varepsilon
\label{eq:thrust}
\end{equation}

where $\eta$ is the conversion efficiency, $\Delta V$ is the vacuum energy released per phase transition, $R/r$ is the torus aspect ratio, and $\varepsilon$ is the asymmetry parameter.

The specific thrust (thrust per input power) scales as:

\begin{equation}
\frac{F}{P} \approx \frac{2\eta\varepsilon}{\omega R} \cdot \frac{R}{r}
\label{eq:specific}
\end{equation}

For a laboratory-scale prototype with parameters $R \approx 0.3$ m, $r \approx 0.1$ m, $\omega \approx 10^6$ s$^{-1}$, $\varepsilon \approx 0.1$, $\eta \approx 0.001$:

\begin{equation}
\frac{F}{P} \approx 6.7 \times 10^{-6} \text{ N/W}
\label{eq:est}
\end{equation}

This is comparable to early ion thruster efficiencies. Optimization of geometric and coupling parameters could yield $10$-$100\times$ improvement.

\section{Stability and Control}

The coupled system exhibits three stability regimes:

\begin{table}[h]
\centering
\begin{tabular}{|c|c|c|}
\hline
\textbf{Regime} & $\varepsilon$ \textbf{range} & \textbf{Behavior} \\
\hline
Subcritical & $\varepsilon < \varepsilon_c$ & No phase transition, elastic response \\
Critical & $\varepsilon = \varepsilon_c \pm \delta$ & Controlled cycling, net impulse \\
Supercritical & $\varepsilon > \varepsilon_c + \delta$ & Runaway, loss of control \\
\hline
\end{tabular}
\caption{Stability regimes of the coupled system.}
\label{tab:stability}
\end{table}

The critical asymmetry is:

\begin{equation}
\varepsilon_c = \frac{V'(\phi_0)}{2\xi B_0^2} \cdot \frac{r}{R}
\label{eq:critical}
\end{equation}

Control is maintained by modulating $B_0$ near the critical value. Each phase transition releases heat $Q_{\text{cycle}} = \Delta V$, requiring thermal management scaling with cycle frequency $\omega_c$.

\section{Falsifiable Predictions}

The PTE framework makes six testable predictions:

\paragraph{P1: Casimir Force Anomalies} At separations $L_k = \phi^k \cdot L_0$ (where $\phi \approx 1.618$ and $L_0 = 2\pi r$), the Casimir force deviates from the ideal parallel-plate value by $>3\%$, arising from the topological vacuum mode structure~\cite{zero2025}.

\paragraph{P2: Anomalous Thrust} RF cavities with asymmetric toroidal dielectric inserts produce thrust at specific resonant frequencies where dilaton coupling is maximal: $F \propto Q \cdot \varepsilon \cdot P_{\text{RF}} / \omega$.

\paragraph{P3: Gamma-Ray Signatures} Vacuum phase transition decay produces monochromatic lines in the $10$-$30$ GeV range, consistent with the excess reported by Totani (2025)~\cite{totani2025}.

\paragraph{P4: Neutron Star Spin-Down} In strong magnetic fields ($\sim 10^{12}$ G), toroidal vacuum drag produces anomalous spin-down torque not explained by magnetic dipole radiation alone.

\paragraph{P5: Dilaton Modulation} The 95.4 GeV diphoton excess shows directional modulation correlated with the local galactic magnetic field.

\paragraph{P6: Apparent Mass Variation} Objects within the toroidal null zone exhibit apparent mass variation proportional to the vacuum stiffness gradient, detectable via precision gravimetry.

\section{Experimental Architecture}

A minimal PTE prototype requires:
\begin{enumerate}
    \item Toroidal RF cavity with high-Q dielectric resonator ($R/r \approx 3$)
    \item Asymmetric field gradient via off-axis dielectric inserts ($\varepsilon \approx 0.05$-$0.2$)
    \item Pulsed RF source ($1$-$10$ kW, resonance-tuned)
    \item Precision thrust stand ($1$ $\mu$N resolution, torsion balance)
\end{enumerate}

The measurement protocol includes:
\begin{itemize}
    \item Thrust vs.\ RF power (tests P2)
    \item Thrust vs.\ asymmetry $\varepsilon$ (tests control model)
    \item Null tests: symmetric configuration and below-threshold power should yield $F = 0$
\end{itemize}

\section{Conclusion}

The Phase-Toroidal Engine offers a theoretically grounded path to propellantless propulsion that avoids the negative energy requirements of earlier proposals. By coupling toroidal field geometry to vacuum phase transition dynamics, the PTE creates directed impulse through topological symmetry breaking. Six falsifiable predictions are identified, testable with existing or near-term experimental infrastructure.

\begin{thebibliography}{9}

\bibitem{weinstock2026}
J.~Weinstock, ``Toroidal Phase-Transition Model (TPTM),'' Zenodo 18138596 (2026).

\bibitem{caputo2025}
R.~N.~Caputo, ``Photonic-Conjugated Model (PCM),'' Zenodo 17823666 (2025).

\bibitem{zero2025}
Zero 4All, ``Toroidal Torsion Field Theory (TTFT),'' Zenodo 17533268 (2025).

\bibitem{hofseth2026}
J.~D.~Hofseth, ``Refractive Vacuum Gravity (RVG),'' Zenodo 18663961 (2026).

\bibitem{blume2025}
``Blume-Capel degrees of freedom in toroidal spin ice,'' \arXiv{2509.01522} (2025).

\bibitem{totani2025}
T.~Totani, ``20 GeV gamma-ray excess from the Galactic Center,'' (2025).

\bibitem{woodward2026}
V.~Woodward, ``Toroidal Scale Dynamics (TSD),'' Zenodo 18357523 (2026).

\bibitem{hofseth2026b}
J.~D.~Hofseth, ``Unified Field Scalar-Hydraulic Drive,'' Zenodo 18652906 (2026).

\end{thebibliography}

\end{document}
